\documentclass{article}
\usepackage[T1]{fontenc}
\usepackage{lmodern}
\usepackage{graphicx}
\usepackage{amsmath,amssymb}
\usepackage[a4paper,margin=2cm]{geometry}
\usepackage[absolute,overlay]{textpos}
\setlength{\TPHorizModule}{1mm}
\setlength{\TPVertModule}{1mm}
\usepackage[indonesian]{babel}
\usepackage{hyperref}
\setlength{\emergencystretch}{2em}
% unit: Halaman 13

\begin{document}

% === Halaman 13 (python/native) ===

2. Blind steganalysis

• Teknik steganalisis yang bekerja pada sembarang algoritma steganografi dan 

sembarang format media.

• Teknik ini mempelajari perbedaan antara statistik cover-object dan stego-object 

dan membedakannya. Proses pembelajaran (learning) dilakukan dengan melatih 

(training) mesin pada sekumpulan database media. Model machine learning yang 

digunakan misalnya jaringan syaraf tiruan.

• Hasil steganalisis kurang akurat dibandingkan dengan teknik targeted steganalysis, 

tetapi kelebihannya adalah dapat diperluas untuk algoritma yang lain.

13

\end{document}
